% Scope: Develop BOLD signal formation and acquisition trade-offs before a limited statistical treatment; distinguish vascular inference from direct neuronal measurement.
% Status: architecture only; no chapter prose or completed problems yet.
% Prerequisites: Chapters 8, 9 and 16.
% Planned worked examples: TE dependence, temporal sampling and physiological-noise limits.
% Planned figures: Hemodynamic response and temporal sampling.
% Planned challenge problems: Optimize an EPI protocol under dropout constraints; diagnose a confounded activation.
\chapter{Functional MRI}
\label{ch:21}

\section{BOLD contrast physics}
\label{sec:21:bold-contrast-physics}
\subsection{Deoxyhemoglobin and susceptibility}
\subsection{Intravascular and extravascular contributions}
\subsection{Gradient-echo and spin-echo sensitivity}

\section{Neurovascular coupling}
\label{sec:21:neurovascular-coupling}
\subsection{Blood flow, volume and oxygen metabolism}
\subsection{Hemodynamic response and temporal variability}
\subsection{Vascular specificity and limitations}

\section{Acquisition design}
\label{sec:21:acquisition-design}
\subsection{EPI, TE and echo spacing}
\subsection{Spatial resolution, temporal resolution and SMS}
\subsection{Distortion, dropout and physiological noise}

\section{Time-series inference}
\label{sec:21:time-series-inference}
\subsection{Drift, motion and nuisance regressors}
\subsection{General linear model and correlated noise}
\subsection{Multiple comparisons and statistical interpretation}

\section{Limits and validation}
\label{sec:21:limits-and-validation}
\subsection{Confounding and causal claims}
\subsection{Task and resting-state measurements}
\subsection{Reproducibility and acquisition-dependent bias}

\section{Worked examples}
\label{sec:21:worked-examples}
% Planned calculations: TE dependence, temporal sampling and physiological-noise limits.
\subsection{Derivation and quantitative calculation}
\subsection{Assumptions, limiting cases and scanner interpretation}

\section{Challenge problems}
\label{sec:21:challenge-problems}
% Problem design: Optimize an EPI protocol under dropout constraints; diagnose a confounded activation.
% Use challengeproblem and stable labels prob:21:<topic>.
% Complete solutions belong in Appendix E, section for this chapter.
% Use \begin{solution}{prob:21:<topic>} ... \end{solution} there.
\subsection{Derivation and quantitative design}
\subsection{Model criticism and integrated reasoning}
